\title[La sabiduría]{\bfseries{La multiforme sabiduría de Dios}}
\subtitle{``\SENOR{} con sabiduría fundó la tierra'' (Pr 3:19)}
\author{}\institute{}\date{}
\titlegraphic{\includegraphics[width=0.80\linewidth,height=0.36\textheight,keepaspectratio]{portada.jpg}\\[-1pt]
  {\fontsize{5}{6}\selectfont\color{gris}\textit{Nota.} Elaboración propia: Pr 3:19 en texto hebreo consonántico (BHS).}}
\AtBeginSection[]{\begin{frame}\frametitle{Contenidos}\tableofcontents[currentsection]\end{frame}}
\begin{document}
\frame{\titlepage}

\begin{frame}\frametitle{Contenidos}\tableofcontents\end{frame}

% ------------------------------------------------------------------
\section{Contexto}

\begin{frame}{El verso de la obra empieza aquí}
\begin{cita}Así dijo \SENOR: No se alabe el sabio en su sabiduría, ni en su valentía se alabe el valiente, ni el rico se alabe en sus riquezas. Mas alábese en esto el que se hubiere de alabar: en entenderme y conocerme, que yo soy \SENOR, que hago misericordia, juicio y justicia en la tierra; porque estas cosas quiero, dice \SENOR.\fuente{(Jer 9:23--24 [9:22--23], RVR1960)}\end{cita}
\begin{center}\heb{אַל־יִתְהַלֵּל חָכָם בְּחָכְמָתוֹ}\quad$\longleftrightarrow$\quad\heb{הַשְׂכֵּל וְיָדֹעַ אוֹתִי}\end{center}
\small\RaggedRight La primera gloria que Jeremías descarta es la del \destacar{sabio} (\emph{jajam}) en su \destacar{sabiduría} (\emph{jojmá}). La única gloria legítima es \emph{haskel veyadoa otí}: ``entender y conocer'' a Dios. La sabiduría humana se rinde ante la sabiduría de Dios.
\end{frame}

\begin{frame}{¿Dónde se hallará la sabiduría?}
\begin{columns}[t,onlytextwidth]
\begin{column}{0.50\textwidth}
\begin{cita}Mas ¿dónde se hallará la sabiduría? ¿Dónde está el lugar de la inteligencia?\fuente{(Job 28:12, RVR1960)}\end{cita}
\begin{cita}Dios entiende el camino de ella, y conoce su lugar. Porque él mira hasta los fines de la tierra, y ve cuanto hay bajo los cielos. Al dar peso al viento, y poner las aguas por medida\ldots{} Y dijo al hombre: He aquí que el temor del Señor es la sabiduría, y el apartarse del mal, la inteligencia.\fuente{(Job 28:23--25, 28, RVR1960)}\end{cita}
\end{column}
\begin{column}{0.48\textwidth}\small
\begin{itemize}
\item Job 28 recorre las minas de la tierra: el hombre saca oro y piedras preciosas de lo profundo, pero no encuentra la sabiduría.
\item Solo Dios sabe ``su lugar'', porque la usó al pesar el viento y medir las aguas (vv. 25--26).
\item Al ser humano le toca la sabiduría como \destacar{temor del Señor} (v. 28): una relación, no una posesión.
\end{itemize}
\end{column}
\end{columns}
\end{frame}

\begin{frame}{Un atributo de la relación}
\centering
\begin{tikzpicture}[font=\scriptsize]
 \node[caja,text width=3.3cm,align=center] (s) at (-4.2,0) {\textbf{Ser}\\ ¿qué es?\\ Uno · Amor · Espíritu};
 \node[caja,text width=3.3cm,align=center] (i) at (0,0) {\textbf{Identidad}\\ ¿quién es?\\ Santo · Eterno · Soberano · Justo · Inmutable · Omnipotente · Omnipresente · Omnisciente};
 \node[draw=dorado,line width=1.6pt,fill=azul,text=papiro,rounded corners=3pt,text width=3.3cm,align=center,inner sep=5pt] (r) at (4.2,0) {\textbf{Relación}\\ ¿cómo se revela?\\ \textcolor{dorado}{\textbf{Sabiduría}} · Verdadero · Misericordioso · Fiel · Celo · Ira · Bondad · Paciencia · Gracia};
\end{tikzpicture}\par\medskip
\begin{columns}[t,onlytextwidth]
\begin{column}{0.49\textwidth}\small
\destacar{Incomunicable en su plenitud:} ``al único y sabio Dios'' (Ro 16:27, \grc{μόνῳ σοφῷ θεῷ}). Nadie le enseñó ni lo aconsejó (Is 40:13--14).
\end{column}
\begin{column}{0.49\textwidth}\small
\destacar{Comunicable como don:} ``\SENOR{} da la sabiduría'' (Pr 2:6); ``si alguno de vosotros tiene falta de sabiduría, pídala a Dios'' (Stg 1:5).
\end{column}
\end{columns}
\end{frame}

% ------------------------------------------------------------------
\section{Proverbios 8}

\begin{frame}{Texto ancla: Pr 8:22--31}
\fotolado{panel_proverbios8.jpg}{Elaboración propia: Pr 8:30 en texto hebreo consonántico (BHS).}
\begin{columns}[c,onlytextwidth]
\begin{column}{0.58\textwidth}
\begin{cita}[fontupper=\scriptsize\itshape]\SENOR{} me poseía en el principio, ya de antiguo, antes de sus obras. Eternamente tuve el principado, desde el principio, antes de la tierra\ldots{} Cuando formaba los cielos, allí estaba yo; cuando trazaba el círculo sobre la faz del abismo\ldots{} cuando ponía al mar su estatuto, para que las aguas no traspasasen su mandamiento; cuando establecía los fundamentos de la tierra, con él estaba yo ordenándolo todo, y era su delicia de día en día, teniendo solaz delante de él en todo tiempo. Me regocijo en la parte habitable de su tierra; y mis delicias son con los hijos de los hombres.\fuente{(Pr 8:22--23, 27, 29--31, RVR1960)}\end{cita}
\end{column}
\begin{column}{0.40\textwidth}\end{column}
\end{columns}
\end{frame}

\begin{frame}{La sabiduría ``ordenándolo todo''}
\begin{center}\small\heb{וָאֶהְיֶה אֶצְלוֹ אָמוֹן וָאֶהְיֶה שַׁעֲשֻׁעִים יוֹם יוֹם}\ \heb{מְשַׂחֶקֶת לְפָנָיו בְּכָל־עֵת}\end{center}
\scriptsize\renewcommand{\arraystretch}{1.15}
\begin{tabular}{@{}>{\RaggedRight\arraybackslash}p{2.6cm}>{\RaggedRight\arraybackslash}p{1.1cm}>{\RaggedRight\arraybackslash}p{8.1cm}@{}}
\toprule
\textbf{Término} & \textbf{Raíz} & \textbf{Sentido} \\ \midrule
\heb{קָנָנִי} \emph{qanani} (v. 22) & \heb{קנה} & ``me adquirió, me poseyó''. La LXX traduce \grc{ἔκτισέν με}, ``me creó'': una lectura que pesará en los debates cristológicos posteriores. \\
\heb{נִסַּכְתִּי} \emph{nissajti} (v. 23) & \heb{נסך} & ``fui establecida'' (RVR1960: ``tuve el principado''). \\
\heb{חוּג} \emph{jug} (v. 27) & \heb{חוג} & ``círculo, bóveda'': Dios traza con compás el horizonte sobre el abismo. \\
\heb{חֻקּוֹ} \emph{juqqó} (v. 29) & \heb{חקק} & ``su estatuto'': el mar tiene un límite grabado. \\
\heb{אָמוֹן} \emph{amón} (v. 30) & \heb{אמן} & ``artesano, maestro de obras'' (cf. \heb{אָמָּן} \emph{ommán}, Cnt 7:1 [7:2]); otros leen ``niño criado'' (\heb{אָמוּן}). RVR1960: ``ordenándolo todo''. LXX: \grc{ἁρμόζουσα}, ``ajustando, armonizando''. \\
\heb{מְשַׂחֶקֶת} \emph{mesajéqet} (vv. 30--31) & \heb{שׂחק} & ``jugando, regocijándose'': la creación es obra de deleite, no de esfuerzo. \\
\bottomrule
\end{tabular}\par\smallskip
\precision{la Sabiduría de Pr 8 es una personificación poética. Leerla como figura de Cristo es lectura tipológica (ver ``En Cristo''); el texto mismo no la identifica con una persona divina.}
\end{frame}

\begin{frame}{``Con sabiduría fundó la tierra''}
\begin{columns}[t,onlytextwidth]
\begin{column}{0.54\textwidth}
\begin{center}\small\heb{יְהוָה בְּחָכְמָה יָסַד־אָרֶץ כּוֹנֵן שָׁמַיִם בִּתְבוּנָה}\\[1pt]\heb{בְּדַעְתּוֹ תְּהוֹמוֹת נִבְקָעוּ}\end{center}
\begin{cita}[fontupper=\scriptsize\itshape]\SENOR{} con sabiduría fundó la tierra; afirmó los cielos con inteligencia. Con su ciencia los abismos fueron divididos, y destilan rocío los cielos.\fuente{(Pr 3:19--20, RVR1960)}\end{cita}
\begin{center}\small\heb{עֹשֵׂה אֶרֶץ בְּכֹחוֹ מֵכִין תֵּבֵל בְּחָכְמָתוֹ}\end{center}
\begin{cita}[fontupper=\scriptsize\itshape]El que hizo la tierra con su poder, el que puso en orden el mundo con su saber, y extendió los cielos con su sabiduría.\fuente{(Jer 10:12, RVR1960)}\end{cita}
\end{column}
\begin{column}{0.44\textwidth}
\centering
\begin{tikzpicture}[font=\scriptsize]
 \node[draw=dorado,line width=1.4pt,fill=azul,text=papiro,rounded corners=3pt,text width=2.4cm,align=center,inner sep=3pt] (d) at (0,0) {\textbf{\SENOR}, el Arquitecto};
 \node[caja,text width=1.6cm,align=center,inner sep=3pt] (a) at (-1.85,-1.25) {\heb{חָכְמָה}\\ funda la tierra};
 \node[caja,text width=1.6cm,align=center,inner sep=3pt] (c) at (0,-1.25) {\heb{דַּעַת}\\ divide los abismos};
 \node[caja,text width=1.6cm,align=center,inner sep=3pt] (b) at (1.85,-1.25) {\heb{תְּבוּנָה}\\ afirma los cielos};
 \foreach \n in {a,b,c}\draw[-{Latex[length=2mm]},dorado,line width=1pt] (d)--(\n);
\end{tikzpicture}\par\smallskip
\scriptsize\RaggedRight La misma tríada vuelve en Bezalel (Éx 31:3): la sabiduría que hizo el cosmos es la que Dios da al artesano del tabernáculo.\par\smallskip
\heb{מֵכִין} (\emph{mejín}, de \heb{כון}): ``establecer, dejar a punto''. El mundo no solo fue hecho: fue \destacar{ajustado}.\par\smallskip
\precision{en Jer 10:12 la RVR1960 vierte \emph{jojmá} como ``saber'' y \emph{tevuná} como ``sabiduría''.}
\end{column}
\end{columns}
\end{frame}

% ------------------------------------------------------------------
\section{La palabra}

\begin{frame}{La raíz \heb{ח־כ־ם}}
\begin{columns}[c,onlytextwidth]
\begin{column}{0.44\textwidth}
\centering
\begin{tikzpicture}[font=\scriptsize]
 \node[draw=dorado,line width=1.6pt,fill=azul,text=papiro,rounded corners=3pt,text width=2.6cm,align=center,inner sep=6pt,font=\small] (r) at (0,0) {\heb{חָכְמָה}\\ \textbf{\emph{jojmá}}};
 \node[caja,text width=2.4cm,align=center] (a) at (-1.6,1.9) {habilidad experta};
 \node[caja,text width=2.4cm,align=center] (b) at (1.6,1.9) {destreza para hacer algo bien};
 \node[caja,text width=2.6cm,align=center] (c) at (0,-1.9) {capacidad de producir un resultado bien hecho};
 \foreach \n in {a,b,c}\draw[dorado,line width=1.1pt](r)--(\n);
\end{tikzpicture}
\end{column}
\begin{column}{0.54\textwidth}\small
\begin{itemize}
\item Raíz \heb{ח־כ־ם} (\emph{j-k-m}; Strong H2449--H2451). En hebreo bíblico no designa ideas abstractas, sino una \destacar{capacidad práctica}.
\item Por eso un artesano, un gobernante o un justo pueden tener \emph{jojmá}.
\item Siempre implica \destacar{capacidad eficaz con propósito correcto}: saber hacer, y hacerlo bien.
\item En arameo bíblico: \heb{חָכְמָה} \emph{jojmá} (Dn 2:20--23); en la LXX y en el NT: \grc{σοφία} (\emph{sophía}).
\end{itemize}
\end{column}
\end{columns}
\end{frame}

\begin{frame}{Los ámbitos de la \emph{jojmá}}
\scriptsize\renewcommand{\arraystretch}{1.2}
\begin{tabular}{@{}>{\RaggedRight\arraybackslash}p{2.0cm}>{\RaggedRight\arraybackslash}p{1.9cm}>{\RaggedRight\arraybackslash}p{4.6cm}>{\RaggedRight\arraybackslash}p{3.2cm}@{}}
\toprule
\textbf{Ámbito} & \textbf{Ejemplo} & \textbf{Descripción} & \textbf{Texto} \\ \midrule
Destreza, técnica & Bezalel & hacer correctamente una obra (arte, oficio) bajo dirección divina & Éx 31:3: \heb{בְּחָכְמָה וּבִתְבוּנָה וּבְדַעַת} \\
Gobierno, juicio & Salomón & discernimiento para gobernar y juzgar con justicia & 1~R 3:9, 28: \heb{לֵב שֹׁמֵעַ}, ``corazón que escucha'' \\
Consejo, dirección & Jetro & ver un problema y dar una solución sabia y ordenada & Éx 18:17--21: \heb{אִיעָצְךָ}, ``te aconsejaré'' \\
Vida justa & el justo (\emph{tsaddiq}) & conducta recta que refleja sabiduría en la vida diaria & Pr 20:7 \\
Hablar sabio & el justo & expresar sabiduría con la boca: enseñanza, consejo & Sal 37:30: \heb{פִּי־צַדִּיק יֶהְגֶּה חָכְמָה} \\
Consejo colectivo & los consejeros & seguridad y dirección por muchos consejos & Pr 11:14; 20:18 \\
\bottomrule
\end{tabular}\par\medskip
\begin{cita}y lo he llenado del Espíritu de Dios, en sabiduría y en inteligencia, en ciencia y en todo arte, para inventar diseños, para trabajar en oro, en plata y en bronce.\fuente{(Éx 31:3--4, RVR1960)}\end{cita}
\end{frame}

\begin{frame}{``Sabiduría'' en la RVR1960}
\begin{columns}[c,onlytextwidth]
\begin{column}{0.40\textwidth}\footnotesize
\begin{itemize}
\item Tu material: \destacar{226 veces}.
\item Conteo propio, sin los títulos de sección: ``sabiduría(s)'' \textbf{245 veces} en 233 versículos (195 AT, 50 NT).
\item Con los títulos que agregan los editores sube a 262.
\item Dos de cada tres veces traduce \emph{jojmá}; en el NT, siempre \grc{σοφία}.
\end{itemize}
\end{column}
\begin{column}{0.58\textwidth}
\scriptsize\renewcommand{\arraystretch}{1.12}
\begin{tabular}{@{}>{\RaggedRight\arraybackslash}p{3.4cm}>{\RaggedRight\arraybackslash}p{2.4cm}r@{}}
\toprule
\textbf{Término} & \textbf{Sentido} & \textbf{Veces} \\ \midrule
\heb{חָכְמָה} \emph{jojmá} (y arameo \emph{jojmá}, 7) & destreza sabia & 161 \\
\heb{דַּעַת} \emph{da'at}, \heb{דֵּעָה} \emph{de'á} & conocimiento & 24 \\
\heb{שֵׂכֶל} \emph{sékel}; \heb{הִשְׂכִּיל} \emph{hiskil} & perspicacia & 4 \\
\heb{תּוּשִׁיָּה} \emph{tushiyá} & sabiduría eficaz & 3 \\
\heb{בִּינָה} \emph{biná} & discernimiento & 1 \\
\heb{כֹּחַ לֵב} \emph{kóaj lev} (Job 36:5) & ``fuerza de corazón'' & 1 \\
sin término equivalente (Job 40:2) & --- & 1 \\
\grc{σοφία} \emph{sophía} & sabiduría & 50 \\
\bottomrule
\end{tabular}
\end{column}
\end{columns}\par\smallskip
\precision{conteo propio con la RVR1960 y la BHS con números Strong; los casos sin etiqueta se asignaron por la raíz del verso hebreo. Conviene verificarlo en e-Sword: la diferencia con 226 puede venir de contar versículos o de otra edición.}
\end{frame}

\begin{frame}{La familia de la sabiduría}
\scriptsize\renewcommand{\arraystretch}{1.16}
\begin{tabular}{@{}>{\RaggedRight\arraybackslash}p{2.3cm}>{\RaggedRight\arraybackslash}p{1.0cm}>{\RaggedRight\arraybackslash}p{5.0cm}>{\RaggedRight\arraybackslash}p{3.4cm}@{}}
\toprule
\textbf{Término} & \textbf{Raíz} & \textbf{Imagen y sentido} & \textbf{Textos} \\ \midrule
\heb{חָכְמָה} \emph{jojmá} & \heb{חכם} & destreza para hacer bien las cosas; el ``arte'' de vivir y de obrar & Pr 3:19; Sal 104:24 \\
\heb{תְּבוּנָה} \emph{tevuná}, \heb{בִּינָה} \emph{biná} & \heb{בין} & ``entre'': distinguir una cosa de otra, ver la diferencia & Pr 3:19; Is 40:28; 1~R 3:9 \\
\heb{דַּעַת} \emph{da'at} & \heb{ידע} & conocimiento por experiencia y trato personal & Pr 3:20; Sal 19:2 [19:3] \\
\heb{עֵצָה} \emph{etsá} & \heb{יעץ} & consejo, plan, designio & Job 12:13; Is 28:29 \\
\heb{תּוּשִׁיָּה} \emph{tushiyá} & \heb{ישׁה}? & sabiduría que da resultado; éxito sólido & Is 28:29; Pr 2:7 \\
\heb{מַחֲשָׁבָה} \emph{majashavá} & \heb{חשׁב} & ``diseño, plan'' (Bezalel) y ``pensamiento'' de Dios & Éx 31:4; Is 55:8 \\
\bottomrule
\end{tabular}\par\medskip
\begin{columns}[t,onlytextwidth]
\begin{column}{0.49\textwidth}
\begin{cita}Con Dios está la sabiduría y el poder; suyo es el consejo y la inteligencia.\fuente{(Job 12:13, RVR1960)}\end{cita}
\end{column}
\begin{column}{0.49\textwidth}
\begin{cita}para hacer maravilloso el consejo y engrandecer la sabiduría.\fuente{(Is 28:29, RVR1960)}\end{cita}
\end{column}
\end{columns}
\precision{la etimología de \emph{tushiyá} es discutida (HALOT).}
\end{frame}

\begin{frame}{Del hebreo a Pablo}
\centering
\begin{tikzpicture}[font=\scriptsize,node distance=0.75cm]
 \node[caja,text width=3.3cm,align=center] (h) {\textbf{Hebreo}\\ \heb{חָכְמָה} \emph{jojmá}\\ destreza del artesano y del Creador\\ {\tiny\color{gris}(Pr 3:19; Éx 31:3)}};
 \node[caja,text width=3.3cm,align=center,right=of h] (g) {\textbf{LXX}\\ \grc{σοφία} \emph{sophía}\\ ``Dios con la sabiduría fundó la tierra''\\ {\tiny\color{gris}(Pr 3:19 LXX)}};
 \node[draw=dorado,line width=1.4pt,fill=azul,text=papiro,rounded corners=3pt,text width=3.3cm,align=center,inner sep=5pt,right=of g] (p) {\textbf{Pablo}\\ \grc{θεοῦ σοφία}\\ ``Cristo, poder de Dios y sabiduría de Dios''\\ {\tiny(1~Co 1:24)}};
 \draw[-{Latex[length=3mm]},dorado,line width=1.4pt] (h)--(g); \draw[-{Latex[length=3mm]},dorado,line width=1.4pt] (g)--(p);
\end{tikzpicture}\par\medskip
{\small\itshape\color{gris}La palabra que en Proverbios describe cómo Dios hizo el mundo, Pablo la usa para describir a Cristo crucificado.}\par\medskip
\begin{cita}Porque los judíos piden señales, y los griegos buscan sabiduría; pero nosotros predicamos a Cristo crucificado\ldots{} mas para los llamados, así judíos como griegos, Cristo poder de Dios, y sabiduría de Dios.\fuente{(1~Co 1:22--24, RVR1960)}\end{cita}
\end{frame}

% ------------------------------------------------------------------
\section{Definición}

\begin{frame}{Qué es la sabiduría según la Escritura}
\centering
\begin{tikzpicture}
 \foreach \x/\n/\t/\d/\r in {0/1/{Maestría creadora}/{diseña y funda el mundo con arte}/{Pr 3:19; 8:30 (Salomón)}, 4.2/2/{Gobierno de los tiempos}/{sostiene la creación y conduce la historia}/{Dn 2:20--21 (Daniel)}, 8.4/3/{Consejo y redención}/{diseño perfecto que nadie le enseñó}/{Is 40:13--14 (Isaías)}}{
  \node[caja,anchor=north] (b\n) at (\x,0) {\parbox[t][2.2cm][t]{3.55cm}{\RaggedRight
    {\small\textbf{\textcolor{azul}{\t}}}\\[5pt]\small \d\par\vfill{\scriptsize\color{gris}(\r)}}};
  \node[circle,fill=azul,text=papiro,font=\scriptsize\bfseries,inner sep=1.6pt] at (b\n.north west) {\n};}
\end{tikzpicture}\par\medskip
\begin{columns}[t,onlytextwidth]
\begin{column}{0.49\textwidth}
\begin{cita}Sea bendito el nombre de Dios de siglos en siglos, porque suyos son el poder y la sabiduría. El muda los tiempos y las edades; quita reyes, y pone reyes; da la sabiduría a los sabios.\fuente{(Dn 2:20--21, RVR1960)}\end{cita}
\end{column}
\begin{column}{0.49\textwidth}
\begin{cita}¿Quién enseñó al Espíritu de \SENOR, o le aconsejó enseñándole? ¿A quién pidió consejo para ser avisado?\fuente{(Is 40:13--14, RVR1960)}\end{cita}
\end{column}
\end{columns}
\small\centering No es solo conocimiento: es \destacar{maestría creadora}, \destacar{habilidad para sostener} su creación y sus tiempos y \destacar{diseño perfecto} de redención y consejo.
\end{frame}

\begin{frame}{Aplicado a Dios}
\begin{columns}[c,onlytextwidth]
\begin{column}{0.50\textwidth}
\small\RaggedRight La sabiduría de Dios es:
\begin{enumerate}
\item el \destacar{``arte''} con el que estructuró el universo;
\item el \destacar{orden} detrás de las leyes de la naturaleza;
\item la \destacar{coherencia} que sostiene y guía la historia;
\item el \destacar{diseño previo} a que existieran tiempo, espacio y materia;
\item el diseño perfecto de la \destacar{redención}.
\end{enumerate}
\end{column}
\begin{column}{0.48\textwidth}
\begin{tcolorbox}[colback=papiro,colframe=azul,boxrule=1pt,arc=3pt]
\small La \heb{חָכְמָה} (\emph{jojmá}) de Dios es su capacidad perfecta y activa para obrar con maestría, mediante la cual \destacar{diseña, ordena, sostiene y dirige} la creación y la historia conforme a su propósito y para su gloria.
\end{tcolorbox}
\end{column}
\end{columns}
\end{frame}

\begin{frame}{Definición desarrollada}
\begin{tcolorbox}[colback=papiro,colframe=azul,boxrule=1pt,arc=3pt]
\small La \heb{חָכְמָה} (\emph{jojmá}) de Dios es un atributo intrínseco e inseparable de su naturaleza, que consiste en su capacidad perfecta y activa de obrar con maestría, mediante la cual Él ---en la unidad absoluta de su ser--- \destacar{sincroniza} de manera perfecta su omnisciencia y omnipresencia con su poder, amor, misericordia y justicia, haciendo que todos sus atributos operen en perfecta armonía para ordenar, sostener y dirigir la creación, la providencia y la redención, cumpliendo su propósito eterno, manifestando su gloria y obrando el bien óptimo de sus hijos.
\end{tcolorbox}
\begin{columns}[t,onlytextwidth]
\begin{column}{0.49\textwidth}\scriptsize
``En la unidad absoluta de su ser'': en Dios los atributos no son piezas separadas que haya que ajustar; cada uno es Dios mismo. Hablar de ``sincronizar'' es una analogía de nuestra manera de conocerlo.
\end{column}
\begin{column}{0.49\textwidth}
\precision{la definición es del autor de la serie. La nota sobre la unidad (simplicidad divina) sigue la teología clásica (Berkhof, 1938; Charnock, 1682/1996).}
\end{column}
\end{columns}
\end{frame}

\begin{frame}{Omnisciencia y sabiduría}
\scriptsize\renewcommand{\arraystretch}{1.25}
\begin{tabular}{@{}>{\RaggedRight\arraybackslash}p{2.0cm}>{\RaggedRight\arraybackslash}p{4.9cm}>{\RaggedRight\arraybackslash}p{4.9cm}@{}}
\toprule
 & \textbf{Omnisciencia} & \textbf{Sabiduría} \\ \midrule
Definición & Dios conoce todas las cosas: pasadas, presentes y futuras. Posee conocimiento perfecto de toda la realidad. & Capacidad divina de usar los mejores medios para lograr los mejores fines, para su gloria. Habilidad perfecta para ordenar todo según su propósito. \\
Pregunta & ¿Qué sabe Dios? & ¿Qué hace Dios con lo que sabe? \\
Imagen & el mapa completo & el camino trazado sobre el mapa \\
Texto & ``conoce lo que está en tinieblas'' (Dn 2:22) & ``Hiciste todas ellas con sabiduría'' (Sal 104:24) \\
\bottomrule
\end{tabular}\par\medskip
\begin{tcolorbox}[colback=papiro,colframe=dorado,boxrule=1.2pt,arc=3pt]
\centering\small Pink: la sabiduría de Dios se muestra ``en la selección de los \destacar{mejores fines posibles} y los \destacar{medios más adecuados} para cumplir sus decretos'' (Pink, 2020, cap. 2).
\end{tcolorbox}
\end{frame}

% ------------------------------------------------------------------
\section{La creación}

\begin{frame}{``En el principio''}
\begin{cita}En el principio creó Dios los cielos y la tierra.\fuente{(Gn 1:1, RVR1960)}\end{cita}
\begin{center}\heb{בְּרֵאשִׁית בָּרָא אֱלֹהִים אֵת הַשָּׁמַיִם וְאֵת הָאָרֶץ}\end{center}
\centering
\begin{tikzpicture}
 \foreach \x/\n/\t/\h/\d in {0/1/{Tiempo}/{\heb{בְּרֵאשִׁית} ``en el principio''}/{pasado, presente y futuro: la historia ocurre dentro de esta dimensión}, 4.2/2/{Espacio}/{\heb{הַשָּׁמַיִם} ``los cielos''}/{largo, ancho y alto: permite movimiento, relación e interacción}, 8.4/3/{Materia}/{\heb{הָאָרֶץ} ``la tierra''}/{sólido, líquido y gaseoso: un universo coherente donde se puede vivir}}{
  \node[caja,anchor=north] (b\n) at (\x,0) {\parbox[t][2.0cm][t]{3.55cm}{\RaggedRight
    {\small\textbf{\textcolor{azul}{\t}}}\enspace{\scriptsize\h}\\[4pt]\scriptsize \d}};
  \node[circle,fill=azul,text=papiro,font=\scriptsize\bfseries,inner sep=1.6pt] at (b\n.north west) {\n};}
\end{tikzpicture}\par\medskip
\precision{leer Gn 1:1 como ``tiempo, espacio y materia'' es una lectura del autor (de raíz creacionista), no un análisis gramatical. ``Los cielos y la tierra'' es en hebreo un merismo: ``todo lo que existe''.}
\end{frame}

\begin{frame}{La creación, un sistema ordenado}
\begin{columns}[t,onlytextwidth]
\begin{column}{0.52\textwidth}
\begin{cita}¡Cuán innumerables son tus obras, oh \SENOR! Hiciste todas ellas con sabiduría; la tierra está llena de tus beneficios.\fuente{(Sal 104:24, RVR1960)}\end{cita}
\begin{center}\heb{כֻּלָּם בְּחָכְמָה עָשִׂיתָ}\end{center}
\scriptsize \heb{קִנְיָנֶךָ} (\emph{qinyaneja}), ``tus posesiones, tus criaturas'' (RVR1960: ``beneficios''), de la misma raíz \heb{קנה} de Pr 8:22: lo que Dios ``adquirió'' y hizo con sabiduría.
\end{column}
\begin{column}{0.46\textwidth}\small
Dios no creó un universo caótico: creó un mundo con \destacar{estructura, leyes y equilibrio}.
\begin{itemize}
\item ciclos de la naturaleza;
\item equilibrio de los ecosistemas;
\item interacción entre las criaturas;
\item dependencia mutua en la vida.
\end{itemize}
Las criaturas se mueven en el tiempo, el espacio y la materia \destacar{en armonía}.
\end{column}
\end{columns}
\end{frame}

\begin{frame}{¿Quién midió las aguas?}
\begin{columns}[c,onlytextwidth]
\begin{column}{0.52\textwidth}
\begin{cita}¿Quién midió las aguas con el hueco de su mano y los cielos con su palmo, con tres dedos juntó el polvo de la tierra, y pesó los montes con balanza y con pesas los collados?\fuente{(Is 40:12, RVR1960)}\end{cita}
\begin{cita}¿No has sabido, no has oído que el Dios eterno es \SENOR, el cual creó los confines de la tierra? No desfallece, ni se fatiga con cansancio, y su entendimiento no hay quien lo alcance.\fuente{(Is 40:28, RVR1960)}\end{cita}
\end{column}
\begin{column}{0.46\textwidth}
\scriptsize\renewcommand{\arraystretch}{1.15}
\begin{tabular}{@{}>{\RaggedRight\arraybackslash}p{1.8cm}>{\RaggedRight\arraybackslash}p{3.3cm}@{}}
\toprule
\heb{מָדַד} \emph{madad} & ``midió'': las medidas de un ingeniero. \\
\heb{תִּכֵּן} \emph{tikkén} & ``ajustó, calibró'' (v. 12, los cielos); el mismo verbo del v. 13: ``¿quién \emph{calibró} el Espíritu del \SENOR?''. \\
\heb{שָׁקַל} \emph{shaqal} & ``pesó'' con balanza. \\
\heb{אֵין חֵקֶר לִתְבוּנָתוֹ} & ``no hay sondeo para su entendimiento'' (v. 28). \\
\bottomrule
\end{tabular}\par\smallskip
\RaggedRight El Creador calibra el universo; a Él nadie lo calibra.
\end{column}
\end{columns}
\end{frame}

\begin{frame}{El ser humano: la culminación}
\begin{columns}[c,onlytextwidth]
\begin{column}{0.50\textwidth}
\begin{cita}Porque tú formaste mis entrañas; tú me hiciste en el vientre de mi madre. Te alabaré; porque formidables, maravillosas son tus obras\ldots{} Mi embrión vieron tus ojos, y en tu libro estaban escritas todas aquellas cosas que fueron luego formadas, sin faltar una de ellas.\fuente{(Sal 139:13--14, 16, RVR1960)}\end{cita}
\end{column}
\begin{column}{0.48\textwidth}
\scriptsize\renewcommand{\arraystretch}{1.15}
\begin{tabular}{@{}>{\RaggedRight\arraybackslash}p{2.1cm}>{\RaggedRight\arraybackslash}p{3.3cm}@{}}
\toprule
\heb{תְּסֻכֵּנִי} \emph{tesukkeni} (v. 13) & ``me tejiste'' (\heb{סכך}). \\
\heb{רֻקַּמְתִּי} \emph{ruqqamti} (v. 15) & ``fui bordado'' (\heb{רקם}): el verbo del bordado del tabernáculo (Éx 35:35). \\
\heb{גָּלְמִי} \emph{golmí} (v. 16) & ``mi embrión'': la masa aún sin forma. \\
\heb{נִפְלָאִים} \emph{niflaím} (v. 14) & ``maravillosos'', obras que superan lo común. \\
\bottomrule
\end{tabular}\par\medskip
\small\RaggedRight La vida humana no es un accidente biológico, sino parte del \destacar{diseño consciente} de Dios.
\end{column}
\end{columns}
\end{frame}

\begin{frame}{Un proceso extraordinariamente coordinado}
\centering
\resizebox{0.86\linewidth}{!}{%
\begin{tikzpicture}[font=\scriptsize]
 % divisiones
 \foreach \i/\n/\x in {0/1/0,1/2/1.7,2/4/3.4,3/8/5.1,4/16/6.8}{
   \node[circle,draw=azul,fill=celeste!25,minimum size=1.15cm,line width=0.8pt] (c\i) at (\x,0) {};
   \pgfmathtruncatemacro{\m}{\n}
   \ifnum\n=1 \fill[azul!70] (\x,0) circle (0.18);\fi
   \ifnum\n=2 \foreach \a in {0,180}{\fill[azul!70] ($(\x,0)+(\a:0.24)$) circle (0.2);}\fi
   \ifnum\n=4 \foreach \a in {45,135,225,315}{\fill[azul!70] ($(\x,0)+(\a:0.26)$) circle (0.15);}\fi
   \ifnum\n=8 \foreach \a in {0,45,...,315}{\fill[azul!70] ($(\x,0)+(\a:0.32)$) circle (0.1);}\fi
   \ifnum\n=16 \foreach \a in {0,22.5,...,337.5}{\fill[azul!70] ($(\x,0)+(\a:0.36)$) circle (0.065);}\fi
   \node[below=2pt] at (c\i.south) {\textbf{\n}};}
 \foreach \i [evaluate=\i as \j using int(\i+1)] in {0,...,3}\draw[-{Latex[length=2mm]},dorado,line width=1pt] (c\i)--(c\j);
 \node[font=\scriptsize\itshape,text=gris] at (3.4,-1.15) {cigoto $\rightarrow$ mitosis: 1, 2, 4, 8, 16\ldots};
 % diferenciacion
 \node[draw=dorado,line width=1.4pt,fill=azul,text=papiro,rounded corners=3pt,text width=2.2cm,align=center] (adn) at (10.2,0) {\textbf{ADN}\\ información que dirige};
 \draw[-{Latex[length=2mm]},dorado,line width=1pt] (c4)--(adn);
 \foreach \t/\y/\k in {neuronas/1.6/a,células sanguíneas/0.55/b,músculo/-0.55/c,hueso/-1.6/d}{\node[caja,text width=2.0cm,align=center,inner sep=2pt] (d\k) at (13.0,\y) {\t}; \draw[dorado,line width=0.8pt](adn)--(d\k);}
\end{tikzpicture}}\par\smallskip
\begin{columns}[t,onlytextwidth]
\begin{column}{0.49\textwidth}\scriptsize
Las células \destacar{saben} cuándo dividirse, cuándo detenerse, en qué tejido convertirse y dónde ubicarse: cerebro, corazón, pulmones, piel, uñas, cabello.
\end{column}
\begin{column}{0.49\textwidth}\scriptsize
La Biblia lee esa coordinación como sabiduría: el tiempo permite la historia, el espacio la interacción, la materia la vida; y dentro de ese sistema, Dios forma al ser humano.
\end{column}
\end{columns}\par\smallskip
\precision{``saben'' es una forma de hablar: la coordinación la regula la información genética y química de las células.}
\end{frame}

% ------------------------------------------------------------------
\section{Providencia}

\begin{frame}{Dios gobierna la historia}
\begin{columns}[t,onlytextwidth]
\begin{column}{0.50\textwidth}
\begin{cita}El muda los tiempos y las edades; quita reyes, y pone reyes; da la sabiduría a los sabios, y la ciencia a los entendidos.\fuente{(Dn 2:21, RVR1960)}\end{cita}
\begin{center}\heb{וְהוּא מְהַשְׁנֵא עִדָּנַיָּא וְזִמְנַיָּא}\end{center}
\scriptsize Arameo: \heb{עִדָּנִין} (\emph{iddanín}) ``tiempos'' y \heb{זִמְנִין} (\emph{zimnín}) ``plazos, momentos fijados''. El Dios sabio es el Señor del calendario de los imperios.
\end{column}
\begin{column}{0.48\textwidth}\small
La \destacar{providencia} es la ejecución del plan de Dios dentro de la historia: la manera en que Dios sostiene, dirige y gobierna los acontecimientos y la vida de los seres humanos para cumplir su propósito.\par\medskip
La historia no avanza al azar. Las naciones pueden parecer poderosas, pero la autoridad última pertenece a Dios, que decide cuándo se levantan y cuándo terminan los reinados.
\end{column}
\end{columns}
\end{frame}

\begin{frame}{Dios dirige el camino de las personas}
\begin{columns}[t,onlytextwidth]
\begin{column}{0.50\textwidth}
\begin{cita}El corazón del hombre piensa su camino; mas \SENOR{} endereza sus pasos.\fuente{(Pr 16:9, RVR1960)}\end{cita}
\begin{center}\heb{לֵב אָדָם יְחַשֵּׁב דַּרְכּוֹ וַיהוָה יָכִין צַעֲדוֹ}\end{center}
\scriptsize \heb{יְחַשֵּׁב} (\emph{yejashshev}, de \heb{חשׁב}): el hombre ``diseña'' su camino. \heb{יָכִין} (\emph{yajín}, de \heb{כון}): Dios ``afirma, establece'' sus pasos, el mismo verbo con que ``puso en orden el mundo con su saber'' (Jer 10:12, \heb{מֵכִין}).
\end{column}
\begin{column}{0.48\textwidth}
\begin{cita}y les ha prefijado el orden de los tiempos, y los límites de su habitación.\fuente{(Hch 17:26, RVR1960)}\end{cita}
\small No elegimos a nuestros padres, ni el país donde nacemos, ni el tiempo histórico en que vivimos. Hacemos planes y tomamos decisiones; la vida ocurre dentro del marco de la providencia, \destacar{sin anular nuestra responsabilidad}.
\end{column}
\end{columns}
\end{frame}

\begin{frame}{``Buena, agradable y perfecta''}
\begin{cita}No os conforméis a este siglo, sino transformaos por medio de la renovación de vuestro entendimiento, para que comprobéis cuál sea la buena voluntad de Dios, agradable y perfecta.\fuente{(Ro 12:2, RVR1960)}\end{cita}
\scriptsize\renewcommand{\arraystretch}{1.2}
\begin{tabular}{@{}>{\RaggedRight\arraybackslash}p{1.4cm}>{\RaggedRight\arraybackslash}p{2.3cm}>{\RaggedRight\arraybackslash}p{4.7cm}>{\RaggedRight\arraybackslash}p{3.4cm}@{}}
\toprule
\textbf{Español} & \textbf{Griego} & \textbf{Acepciones} & \textbf{Idea central} \\ \midrule
Buena & \grc{ἀγαθός} (\emph{agathós}) & bueno, beneficioso, útil, noble; lo que produce bienestar real; lo moralmente correcto y provechoso & lo que produce verdadero bien \\
Agradable & \grc{εὐάρεστος} (\emph{euárestos}) & aceptable, grato, que agrada, aprobado por Dios & lo que agrada plenamente a Dios \\
Perfecta & \grc{τέλειος} (\emph{téleios}) & completo, pleno, maduro, acabado; que alcanza su meta (\grc{τέλος}) & lo que cumple el propósito para el que fue diseñado \\
\bottomrule
\end{tabular}\par\medskip
\small\RaggedRight La providencia no es arbitraria: Dios dirige la historia y nuestras vidas con sabiduría perfecta, hacia un \destacar{\emph{télos}}.
\end{frame}

\begin{frame}{Todas las cosas cooperan}
\begin{columns}[t,onlytextwidth]
\begin{column}{0.52\textwidth}
\begin{cita}Y sabemos que a los que aman a Dios, todas las cosas les ayudan a bien, esto es, a los que conforme a su propósito son llamados. Porque a los que antes conoció, también los predestinó para que fuesen hechos conformes a la imagen de su Hijo.\fuente{(Ro 8:28--29, RVR1960)}\end{cita}
\scriptsize\grc{πάντα συνεργεῖ εἰς ἀγαθόν}: ``todas las cosas \emph{co-operan} (\grc{σύν} + \grc{ἔργον}) para bien''.
\end{column}
\begin{column}{0.46\textwidth}\small
El versículo no dice que todo lo que ocurre es bueno, sino que Dios hace que todo coopere para un bien final. Esto incluye:
\begin{itemize}
\item sufrimiento y pruebas;
\item injusticias y dificultades.
\end{itemize}
El bien final que Dios busca es \destacar{hacernos semejantes a Cristo} (\grc{συμμόρφους τῆς εἰκόνος τοῦ υἱοῦ}, v. 29).
\end{column}
\end{columns}
\precision{algunos manuscritos antiguos (\(\mathfrak{P}^{46}\), A, B) leen \grc{συνεργεῖ ὁ θεός}: ``Dios hace cooperar todas las cosas''. NA28 no incluye \grc{ὁ θεός} en el texto.}
\end{frame}

\begin{frame}{Cuando no entendemos}
\begin{columns}[t,onlytextwidth]
\begin{column}{0.50\textwidth}
\begin{cita}Porque mis pensamientos no son vuestros pensamientos, ni vuestros caminos mis caminos, dijo \SENOR. Como son más altos los cielos que la tierra, así son mis caminos más altos que vuestros caminos.\fuente{(Is 55:8--9, RVR1960)}\end{cita}
\begin{cita}Todo lo hizo hermoso en su tiempo; y ha puesto eternidad en el corazón de ellos, sin que alcance el hombre a entender la obra que ha hecho Dios desde el principio hasta el fin.\fuente{(Ec 3:11, RVR1960)}\end{cita}
\end{column}
\begin{column}{0.48\textwidth}\small
La providencia opera en un nivel de sabiduría mucho más amplio que nuestra perspectiva.\par\medskip
\heb{יָפֶה בְעִתּוֹ} (\emph{yafé be-ittó}): ``hermoso \destacar{en su tiempo}''. Lo que hoy parece fuera de lugar tiene su hora en el diseño completo.\par\medskip
Pink: cuando algo no se puede conciliar con nuestras nociones de lo bueno y lo sabio, debemos ``juzgar el todo por la parte; es decir, lo desconocido por medio de lo conocido'' (Pink, 2020, cap. 2).
\end{column}
\end{columns}
\end{frame}

% ------------------------------------------------------------------
\section{En Israel}

\begin{frame}{Un encuentro: José en Egipto}
\centering
\begin{tikzpicture}[font=\scriptsize]
 \draw[line width=1.6pt,azul,-{Latex[length=3mm]}] (-0.3,0) -- (12.2,0);
 \foreach \x/\t/\b/\r/\pos in {0.4/{Vendido}/{veinte piezas de plata}/{Gn 37:28}/above,
   2.4/{Potifar}/{``\SENOR{} estaba con José''}/{Gn 39:2}/below,
   4.4/{La cárcel}/{acusado injustamente}/{Gn 39:20}/above,
   6.4/{Ante Faraón}/{``no hay entendido ni sabio como tú''}/{Gn 41:39}/below,
   8.4/{Gobernador}/{pan para el mundo en el hambre}/{Gn 41:41}/above,
   10.6/{``Lo encaminó a bien''}/{su familia preservada}/{Gn 50:20}/below}{
  \fill[dorado] (\x,0) circle (4pt);
  \node[\pos=7pt,text width=2.0cm,align=center] at (\x,0) {\textbf{\textcolor{azul}{\t}}\\ \b\\ {\tiny\color{gris}(\r)}};}
 \draw[decorate,decoration={brace,amplitude=5pt,mirror},gris] (0.2,-2.25) -- (4.6,-2.25) node[midway,below=6pt,text=gris] {\emph{lo que parece desorden}};
 \draw[decorate,decoration={brace,amplitude=5pt,mirror},gris] (6.2,-2.25) -- (11.4,-2.25) node[midway,below=6pt,text=gris] {\emph{lo que Dios estaba tejiendo}};
\end{tikzpicture}\par\medskip
\small\RaggedRight Trece años de traición, esclavitud y cárcel. Ninguno de esos pasos parecía conducir al trono; todos conducían a él.
\end{frame}

\begin{frame}{Dios lo encaminó a bien}
\fotolado{panel_genesis50.jpg}{Elaboración propia: Gn 50:20 en texto hebreo consonántico (BHS).}
\begin{columns}[c,onlytextwidth]
\begin{column}{0.58\textwidth}
\begin{cita}Vosotros pensasteis mal contra mí, mas Dios lo encaminó a bien, para hacer lo que vemos hoy, para mantener en vida a mucho pueblo.\fuente{(Gn 50:20, RVR1960)}\end{cita}
\begin{center}\heb{וְאַתֶּם חֲשַׁבְתֶּם עָלַי רָעָה אֱלֹהִים חֲשָׁבָהּ לְטֹבָה}\end{center}
\scriptsize\RaggedRight Un mismo verbo, \heb{חָשַׁב} (\emph{jashav}), ``planear, diseñar'', con dos sujetos: \emph{vosotros} planearon el mal; \emph{Dios} lo planeó para bien. Es el verbo del artesano Bezalel, ``para inventar diseños'' (\heb{לַחְשֹׁב מַחֲשָׁבֹת}, Éx 31:4).\par\smallskip
\small Dos intenciones simultáneas sobre el mismo hecho: la maldad de los hermanos y el \destacar{diseño} de Dios, que transforma el mal humano en instrumento de un bien mayor.
\end{column}
\begin{column}{0.40\textwidth}\end{column}
\end{columns}
\end{frame}

\begin{frame}{José lee su propia historia}
\begin{columns}[t,onlytextwidth]
\begin{column}{0.50\textwidth}
\begin{cita}Ahora, pues, no os entristezcáis, ni os pese de haberme vendido acá; porque para preservación de vida me envió Dios delante de vosotros\ldots{} Así, pues, no me enviasteis acá vosotros, sino Dios.\fuente{(Gn 45:5, 8, RVR1960)}\end{cita}
\begin{cita}Y dijo Faraón a José: Pues que Dios te ha hecho saber todo esto, no hay entendido ni sabio como tú.\fuente{(Gn 41:39, RVR1960)}\end{cita}
\end{column}
\begin{column}{0.48\textwidth}\small
\begin{itemize}
\item \heb{אֵין־נָבוֹן וְחָכָם כָּמוֹךָ}: ``no hay \emph{navón} y \emph{jajam} como tú''. La sabiduría de José es recibida: ``Dios te ha hecho saber''.
\item José no niega el mal de sus hermanos (50:20), pero ve por encima de él la mano de Dios (45:8).
\item Por eso puede perdonar y sustentar: ``yo os sustentaré'' (\heb{אֲכַלְכֵּל}, 50:21).
\end{itemize}
\end{column}
\end{columns}
\end{frame}

% ------------------------------------------------------------------
\section{Pink}

\begin{frame}{Teología: los decretos son sabios}
\begin{columns}[t,onlytextwidth]
\begin{column}{0.55\textwidth}\small
\RaggedRight Pink no dedica un capítulo a la sabiduría; la trata al hablar de los decretos (2020, cap. 2):
\begin{enumerate}
\item Los decretos se llaman ``consejo'' para indicar que son \destacar{perfectamente sabios}; en Dios la sabiduría siempre acompaña a la voluntad (Ef 1:11).
\item Su sabiduría se ve en la elección de los \destacar{mejores fines} y los \destacar{medios más adecuados}.
\item Conocemos una parte muy pequeña de sus obras: hay que juzgar el todo por la parte.
\item Sus decretos son \destacar{libres}: nadie le enseñó ni le aconsejó (Is 40:13--14).
\end{enumerate}
\end{column}
\begin{column}{0.43\textwidth}
\begin{tcolorbox}[colback=papiro,colframe=dorado,boxrule=1.2pt,arc=3pt]
\scriptsize ``Aquel que percibe el admirable ingenio en el funcionamiento de las partes de una máquina\ldots{} es conducido a creer que las demás partes de esta máquina son igualmente admirables.''\par\hfill{\upshape\bfseries(Pink, 2020, cap. 2)}
\end{tcolorbox}
\begin{tcolorbox}[colback=papiro,colframe=dorado,boxrule=1.2pt,arc=3pt]
\scriptsize ``Su justicia es una justicia santa, su sabiduría es una \destacar{sabiduría santa} y su poder es `su santo brazo'.''\par\hfill{\upshape\bfseries(Charnock, citado en Pink, 2020, cap. 8)}
\end{tcolorbox}
\precision{Pink es aquí una referencia; el orden y la estructura del estudio son de la serie.}
\end{column}
\end{columns}
\end{frame}

% ------------------------------------------------------------------
\section{Judaísmo}

\begin{frame}{La Torá, plano del Artesano}
\begin{columns}[t,onlytextwidth]
\begin{column}{0.53\textwidth}\scriptsize
\textbf{\emph{Bereshit Rabá} 1:1} comenta \heb{אָמוֹן} (Pr 8:30) y da cuatro sentidos: ``pedagogo'', ``cubierto'', ``escondido'' y ``grande''. Luego compara a Dios con un rey que construye un palacio: no lo hace de su propia idea, sino con un artesano, que tiene \heb{דִּפְתְּרָאוֹת וּפִנְקְסָאוֹת}, ``rollos y cuadernos'' de planos.
\begin{center}\heb{הַתּוֹרָה אוֹמֶרֶת אֲנִי הָיִיתִי כְּלִי אֻמְנוּתוֹ}\\ \heb{שֶׁל הַקָּדוֹשׁ בָּרוּךְ הוּא}\end{center}
``La Torá dice: yo era el instrumento del oficio del Santo, bendito sea''. Dios miraba en la Torá y creaba el mundo.
\end{column}
\begin{column}{0.45\textwidth}
\centering
\begin{tikzpicture}[font=\scriptsize,baseline=(current bounding box.north)]
 \node[caja,text width=3.2cm,align=center] (a) at (0,1.5) {\textbf{Pr 8:22}\\ ``\SENOR{} me poseía \emph{en el principio} (\heb{רֵאשִׁית})''};
 \node[caja,text width=3.2cm,align=center] (b) at (0,0) {\textbf{Gn 1:1}\\ ``\heb{בְּרֵאשִׁית} creó Dios''};
 \node[draw=dorado,line width=1.4pt,fill=azul,text=papiro,rounded corners=3pt,text width=3.2cm,align=center,inner sep=4pt] (c) at (0,-1.5) {``Con la Torá (= sabiduría) creó Dios''};
 \draw[-{Latex[length=2mm]},dorado,line width=1pt](a)--(b); \draw[-{Latex[length=2mm]},dorado,line width=1pt](b)--(c);
\end{tikzpicture}\par\smallskip
\precision{interpretación rabínica (\emph{Bereshit Rabá} 1:1, Sefaria). El Targum Neofiti de Gn 1:1 dice: ``Desde el principio, \emph{con sabiduría}, el Memra del \SENOR{} creó y perfeccionó los cielos y la tierra'' (McNamara, 1992).}
\end{column}
\end{columns}
\end{frame}

\begin{frame}{Sabiduría, inteligencia y conocimiento}
\begin{columns}[t,onlytextwidth]
\begin{column}{0.50\textwidth}\small
\RaggedRight\textbf{Rashi} (a Éx 31:3) distingue la tríada con que Dios llenó a Bezalel:
\scriptsize\renewcommand{\arraystretch}{1.25}
\begin{tabular}{@{}>{\RaggedRight\arraybackslash}p{1.6cm}>{\RaggedRight\arraybackslash}p{4.2cm}@{}}
\toprule
\heb{חָכְמָה} & \heb{מַה שֶּׁאָדָם שׁוֹמֵעַ מֵאֲחֵרִים וְלָמֵד}: ``lo que uno oye de otros y aprende''. \\
\heb{תְּבוּנָה} & \heb{מֵבִין דָּבָר מִלִּבּוֹ}: ``entiende algo por sí mismo, a partir de lo aprendido''. \\
\heb{דַּעַת} & \heb{רוּחַ הַקֹּדֶשׁ}: ``el Espíritu Santo''. \\
\bottomrule
\end{tabular}
\end{column}
\begin{column}{0.48\textwidth}\small
\textbf{Dos definiciones del sabio:}
\begin{itemize}
\item Ben Zomá (\emph{Pirké Avot} 4:1): \heb{אֵיזֶהוּ חָכָם הַלּוֹמֵד מִכָּל אָדָם}, ``¿quién es sabio? El que aprende de todo hombre'' (Sal 119:99).
\item Los sabios del Néguev a un rey (\emph{Tamid} 32a): \heb{אֵיזֶהוּ חָכָם הָרוֹאֶה אֶת הַנּוֹלָד}, ``¿quién es sabio? El que ve lo que va a nacer'': el que ve el resultado antes de que llegue.
\end{itemize}
\end{column}
\end{columns}\par\smallskip
\precision{lecturas rabínicas (Rashi; Misná; Talmud de Babilonia), citadas de Sefaria. ``Ver lo que va a nacer'' es, en Dios, ver el fin desde el principio (Is 46:10).}
\end{frame}

\begin{frame}{Bendecir la sabiduría de Dios}
\begin{columns}[t,onlytextwidth]
\begin{column}{0.50\textwidth}\small
\RaggedRight\textbf{\emph{Berajot} 58a}: bendiciones al ver sabios.
\begin{itemize}\RaggedRight
\item A los sabios de Israel: \heb{שֶׁחָלַק מֵחׇכְמָתוֹ לִירֵאָיו}, ``que repartió de su sabiduría a los que le temen''.
\item A los sabios de las naciones: \heb{שֶׁנָּתַן מֵחׇכְמָתוֹ לְבָשָׂר וָדָם}, ``que dio de su sabiduría a carne y sangre''.
\item Al ver una multitud de Israel: \heb{חֲכַם הָרָזִים}, ``Sabio de los secretos'', porque ``sus mentes no se parecen unas a otras, ni sus rostros''.
\end{itemize}
\end{column}
\begin{column}{0.48\textwidth}\small
\textbf{\emph{Berajot} 60b}: bendición diaria por el cuerpo (\emph{Asher yatsar}):
\begin{center}\heb{אֲשֶׁר יָצַר אֶת הָאָדָם בְּחׇכְמָה}\\ \heb{וּבָרָא בּוֹ נְקָבִים נְקָבִים חֲלָלִים חֲלָלִים}\end{center}
``\ldots{} que formó al hombre con sabiduría y creó en él orificios y cavidades''; y sigue: si uno de ellos se abriera o se cerrara, no se podría subsistir delante de Dios.\par\smallskip
\scriptsize Es la versión judía de Sal 139:14: el cuerpo funcionando a tiempo es motivo de bendición cada mañana.
\end{column}
\end{columns}
\precision{la oración diaria de la mañana (\emph{Yotser}) recita también Sal 104:24: ``¡Cuán innumerables son tus obras! Todas las hiciste con sabiduría'' (Elbogen, 1993).}
\end{frame}

\begin{frame}{La Sabiduría en el judaísmo del Segundo Templo}
\begin{columns}[t,onlytextwidth]
\begin{column}{0.50\textwidth}\small
\textbf{Sabiduría de Salomón} (s. I a.\,e.\,c., griego):
\begin{center}\grc{ἡ γὰρ πάντων τεχνῖτις ἐδίδαξέ με σοφία}\end{center}
``Porque me enseñó la sabiduría, artesana de todas las cosas'' (Sab 7:22).
\begin{center}\grc{διατείνει δὲ ἀπὸ πέρατος ἐπὶ πέρας εὐρώστως}\\ \grc{καὶ διοικεῖ τὰ πάντα χρηστῶς}\end{center}
``Se extiende con fuerza de un extremo al otro, y gobierna todas las cosas con bondad'' (Sab 8:1).
\end{column}
\begin{column}{0.48\textwidth}\small
\textbf{Ben Sirá} (s. II a.\,e.\,c.):
\begin{center}\grc{Πᾶσα σοφία παρὰ κυρίου}\end{center}
``Toda sabiduría viene del Señor'' (Sir 1:1). En el cap. 24 la Sabiduría busca dónde morar y planta su tienda en Jacob: ``Todo esto es el libro de la alianza del Dios Altísimo, la ley que Moisés nos prescribió'' (Sir 24:23).\par\medskip
\scriptsize \emph{Sabiduría} 7:26 la llama ``resplandor (\grc{ἀπαύγασμα}) de la luz eterna\ldots{} e imagen (\grc{εἰκών}) de su bondad'': vocabulario que reaparece en Col 1:15 y He 1:3.
\end{column}
\end{columns}
\precision{libros deuterocanónicos; traducción propia desde la LXX (Rahlfs-Hanhart, 2006). Se citan como testigos de la fe judía, no como Escritura de esta serie.}
\end{frame}

\begin{frame}{Qumrán: el plan antes de que existieran}
\begin{columns}[t,onlytextwidth]
\begin{column}{0.52\textwidth}\small
\textbf{Regla de la Comunidad} (1QS III 15--17):
\begin{center}\heb{מֵאֵל הַדֵּעוֹת כּוֹל הוֹוֶה וְנִהְיֶה}\\ \heb{וְלִפְנֵי הֱיוֹתָם הֵכִין כּוֹל מַחֲשַׁבְתָּם}\end{center}
``Del Dios de los conocimientos viene todo lo que es y lo que será; y antes de que existieran, \destacar{estableció todo su diseño}; y cuando existen\ldots{} cumplen su obra según su glorioso designio, y nada se puede cambiar. En su mano están las leyes de todo, y él los \destacar{sustenta} en todas sus necesidades'' (\heb{וְהוּאָה יְכַלְכְּלֵם}).
\end{column}
\begin{column}{0.46\textwidth}\small
\textbf{Gran Rollo de los Salmos} (11QSal\textsuperscript{a} XXVI 9--15), ``Himno al Creador'':\par\smallskip
\begin{cita}Bendito el que hace la tierra con su poder, el que establece el mundo con su sabiduría; con su entendimiento extendió los cielos.\fuente{(11QSal\textsuperscript{a} XXVI 13--15)}\end{cita}
\scriptsize Es Jer 10:12 convertido en bendición.\par\smallskip
\precision{1QS según la transcripción de Qumran-Digital; 11QSal\textsuperscript{a}, traducción propia de la edición de Sanders (1965). Nótese \heb{הֵכִין} y \heb{מַחֲשָׁבָה}: las mismas raíces de Jer 10:12 y Gn 50:20.}
\end{column}
\end{columns}
\end{frame}

% ------------------------------------------------------------------
\section{En Cristo}

\begin{frame}{Cristo, sabiduría de Dios}
\begin{columns}[t,onlytextwidth]
\begin{column}{0.52\textwidth}
\begin{cita}[fontupper=\scriptsize\itshape]Pues ya que en la sabiduría de Dios, el mundo no conoció a Dios mediante la sabiduría, agradó a Dios salvar a los creyentes por la locura de la predicación\ldots{} mas para los llamados, así judíos como griegos, Cristo poder de Dios, y sabiduría de Dios. Porque lo insensato de Dios es más sabio que los hombres.\fuente{(1~Co 1:21, 24--25, RVR1960)}\end{cita}
\end{column}
\begin{column}{0.46\textwidth}
\centering
\begin{tikzpicture}[font=\scriptsize,baseline=(current bounding box.north)]
 \node[caja,text width=2.3cm,align=center] (j) at (-1.45,1.4) {\textbf{Judíos}\\ piden señales\\ (\grc{σημεῖα})};
 \node[caja,text width=2.3cm,align=center] (g) at (1.45,1.4) {\textbf{Griegos}\\ buscan sabiduría\\ (\grc{σοφίαν})};
 \node[draw=dorado,line width=1.4pt,fill=azul,text=papiro,rounded corners=3pt,text width=3.6cm,align=center,inner sep=4pt] (c) at (0,-0.4) {\textbf{Cristo crucificado}\\ \grc{θεοῦ δύναμιν καὶ θεοῦ σοφίαν}};
 \draw[-{Latex[length=2mm]},dorado,line width=1pt](j)--(c); \draw[-{Latex[length=2mm]},dorado,line width=1pt](g)--(c);
 \node[text width=4.6cm,align=center,font=\scriptsize\itshape,text=gris] at (0,-1.7) {la respuesta a las dos búsquedas: el poder que pedía Israel y la sabiduría que buscaba Grecia};
\end{tikzpicture}
\end{column}
\end{columns}
\scriptsize\RaggedRight La cruz parece ``locura'' (\grc{μωρία}), pero es la obra maestra del Artesano: en ella se encuentran la justicia que condena y la misericordia que perdona.
\end{frame}

\begin{frame}{``El que se gloría, gloríese en el Señor''}
\begin{cita}Mas por él estáis vosotros en Cristo Jesús, el cual nos ha sido hecho por Dios sabiduría, justificación, santificación y redención; para que, como está escrito: El que se gloría, gloríese en el Señor.\fuente{(1~Co 1:30--31, RVR1960)}\end{cita}
\scriptsize\renewcommand{\arraystretch}{1.2}
\begin{tabular}{@{}>{\RaggedRight\arraybackslash}p{2.0cm}>{\RaggedRight\arraybackslash}p{4.8cm}>{\RaggedRight\arraybackslash}p{4.9cm}@{}}
\toprule
\textbf{Texto} & \textbf{Hebreo / griego} & \textbf{Lectura} \\ \midrule
Jer 9:23 & \heb{אַל־יִתְהַלֵּל חָכָם בְּחָכְמָתוֹ} & ``No se alabe el sabio en su sabiduría'' \\
Jer 9:24 LXX & \grc{ἐν τούτῳ καυχάσθω ὁ καυχώμενος} & ``en esto gloríese el que se gloría'' \\
1~Co 1:31 & \grc{ὁ καυχώμενος ἐν κυρίῳ καυχάσθω} & Pablo resume Jeremías: el sabio deja su sabiduría y se gloría en Cristo, ``hecho sabiduría'' para nosotros \\
\bottomrule
\end{tabular}\par\medskip
\small\RaggedRight El verso clave de esta obra (Jer 9:23--24) aparece cumplido en Pablo: la sabiduría que Jeremías quitaba al sabio, Dios la da en Cristo.
\end{frame}

\begin{frame}{La sabiduría oculta}
\begin{columns}[t,onlytextwidth]
\begin{column}{0.52\textwidth}
\begin{cita}Mas hablamos sabiduría de Dios en misterio, la sabiduría oculta, la cual Dios predestinó antes de los siglos para nuestra gloria, la que ninguno de los príncipes de este siglo conoció; porque si la hubieran conocido, nunca habrían crucificado al Señor de gloria.\fuente{(1~Co 2:7--8, RVR1960)}\end{cita}
\end{column}
\begin{column}{0.46\textwidth}\small
\begin{itemize}
\item \grc{προώρισεν\ldots{} πρὸ τῶν αἰώνων}: ``predestinó antes de los siglos''. El diseño es anterior al tiempo (Pr 8:23).
\item Los poderes de este siglo creyeron destruir a Cristo y cumplieron el plan: Gn 50:20 en grande.
\item \grc{μυστήριον}: no un enigma, sino un plan escondido y ahora revelado (Ro 16:25--26).
\end{itemize}
\end{column}
\end{columns}
\end{frame}

\begin{frame}{``¡Oh profundidad!''}
\fotolado{panel_romanos11.jpg}{Elaboración propia: Ro 11:33 en escritura uncial (texto NA28).}
\begin{columns}[c,onlytextwidth]
\begin{column}{0.58\textwidth}
\begin{cita}¡Oh profundidad de las riquezas de la sabiduría y de la ciencia de Dios! ¡Cuán insondables son sus juicios, e inescrutables sus caminos! Porque ¿quién entendió la mente del Señor? ¿O quién fue su consejero? ¿O quién le dio a él primero, para que le fuese recompensado? Porque de él, y por él, y para él, son todas las cosas. A él sea la gloria por los siglos. Amén.\fuente{(Ro 11:33--36, RVR1960)}\end{cita}
\scriptsize\renewcommand{\arraystretch}{1.12}
\begin{tabular}{@{}>{\RaggedRight\arraybackslash}p{1.2cm}>{\RaggedRight\arraybackslash}p{5.6cm}@{}}
\toprule
Is 40:13 & \heb{מִי־תִכֵּן אֶת־רוּחַ יְהוָה וְאִישׁ עֲצָתוֹ יוֹדִיעֶנּוּ} \\
LXX & \grc{τίς ἔγνω νοῦν κυρίου, καὶ τίς αὐτοῦ σύμβουλος ἐγένετο} \\
Ro 11:34 & \grc{τίς γὰρ ἔγνω νοῦν κυρίου; ἢ τίς σύμβουλος αὐτοῦ ἐγένετο;} \\
\bottomrule
\end{tabular}
\end{column}
\begin{column}{0.40\textwidth}\end{column}
\end{columns}
\end{frame}

\begin{frame}{Después de Israel, la doxología}
\begin{columns}[c,onlytextwidth]
\begin{column}{0.50\textwidth}\small
\begin{itemize}
\item Ro 9--11 recorre la historia de Israel: elección, tropiezo, injerto de los gentiles y la promesa de que ``todo Israel será salvo'' (11:26).
\item Pablo no termina con una explicación, sino con \destacar{adoración}: la sabiduría que ordena la historia supera lo que él puede explicar.
\item Ro 11:34 cita Is 40:13 (LXX); Ro 11:35 recuerda Job 41:11 [41:3]: \heb{מִי הִקְדִּימַנִי וַאֲשַׁלֵּם}.
\end{itemize}
\end{column}
\begin{column}{0.48\textwidth}
\scriptsize\renewcommand{\arraystretch}{1.2}
\begin{tabular}{@{}>{\RaggedRight\arraybackslash}p{1.5cm}>{\RaggedRight\arraybackslash}p{3.9cm}@{}}
\toprule
\grc{βάθος} & ``profundidad'', como un abismo sin fondo. \\
\grc{ἀνεξεραύνητα} & ``imposibles de rastrear'': juicios que no se pueden investigar hasta el fondo. \\
\grc{ἀνεξιχνίαστοι} & ``sin huellas'': caminos de los que no se pueden seguir las pisadas. \\
\grc{ἐξ\ldots{} δι᾽\ldots{} εἰς} & ``de él, por él, para él'': origen, medio y meta de todas las cosas. \\
\bottomrule
\end{tabular}
\end{column}
\end{columns}
\end{frame}

\begin{frame}{La multiforme sabiduría de Dios}
\begin{columns}[t,onlytextwidth]
\begin{column}{0.52\textwidth}
\begin{cita}para que la multiforme sabiduría de Dios sea ahora dada a conocer por medio de la iglesia a los principados y potestades en los lugares celestiales, conforme al propósito eterno que hizo en Cristo Jesús nuestro Señor.\fuente{(Ef 3:10--11, RVR1960)}\end{cita}
\begin{center}\grc{ἡ πολυποίκιλος σοφία τοῦ θεοῦ}\end{center}
\scriptsize \grc{πολυποίκιλος} (\emph{polypoíkilos}): ``de muchos colores, ricamente variada''. \grc{ποικίλος} describe tejidos bordados de varios colores; en la LXX, la túnica de José es \grc{χιτὼν ποικίλος} (Gn 37:3).
\end{column}
\begin{column}{0.46\textwidth}\small
\begin{itemize}
\item El ``misterio'' de Ef 3:6: los gentiles son ``coherederos, miembros del mismo cuerpo'' con Israel.
\item La \destacar{iglesia}, hecha de judíos y gentiles, es la pieza bordada que muestra la sabiduría de Dios a los poderes celestiales.
\item La sabiduría no se exhibe en una idea, sino en un \destacar{pueblo} reconciliado.
\end{itemize}
\precision{la autoría paulina de Efesios y Colosenses se discute en la academia; la misma enseñanza aparece en cartas indiscutidas (Ro 11:33--36; 1~Co 1--2).}
\end{column}
\end{columns}
\end{frame}

\begin{frame}{Todo reunido en Cristo}
\fotolado{panel_colosenses1.jpg}{Elaboración propia: Col 1:17 en escritura uncial (texto NA28).}
\begin{columns}[c,onlytextwidth]
\begin{column}{0.58\textwidth}
\begin{cita}de reunir todas las cosas en Cristo, en la dispensación del cumplimiento de los tiempos, así las que están en los cielos, como las que están en la tierra.\fuente{(Ef 1:10, RVR1960)}\end{cita}
\begin{cita}El es la imagen del Dios invisible, el primogénito de toda creación\ldots{} Y él es antes de todas las cosas, y todas las cosas en él subsisten.\fuente{(Col 1:15, 17, RVR1960)}\end{cita}
\begin{cita}en quien están escondidos todos los tesoros de la sabiduría y del conocimiento.\fuente{(Col 2:3, RVR1960)}\end{cita}
\end{column}
\begin{column}{0.40\textwidth}\end{column}
\end{columns}
\end{frame}

\begin{frame}{Tres palabras con \grc{σύν}}
\centering
\begin{tikzpicture}
 \foreach \x/\n/\g/\t/\d/\r in {0/1/{συνεργεῖ}/{cooperan}/{todas las cosas trabajan juntas para el bien de los que aman a Dios}/{Ro 8:28}, 4.2/2/{συνέστηκεν}/{subsisten}/{todas las cosas se mantienen unidas, cohesionadas, en Cristo}/{Col 1:17}, 8.4/3/{ἀνακεφαλαιώσασθαι}/{reunir bajo una cabeza}/{todas las cosas, del cielo y de la tierra, recapituladas en Cristo}/{Ef 1:10}}{
  \node[caja,anchor=north] (b\n) at (\x,0) {\parbox[t][2.3cm][t]{3.55cm}{\RaggedRight
    {\small\grc{\g}}\\ {\small\textbf{\textcolor{azul}{\t}}}\\[4pt]\scriptsize \d\par\vfill{\scriptsize\color{gris}(\r)}}};
  \node[circle,fill=azul,text=papiro,font=\scriptsize\bfseries,inner sep=1.6pt] at (b\n.north west) {\n};}
\end{tikzpicture}\par\medskip
\begin{columns}[t,onlytextwidth]
\begin{column}{0.55\textwidth}\small
Pablo describe la sabiduría de Dios en Cristo con palabras de \destacar{unión}: lo que coopera en la historia (Ro 8:28), lo que se mantiene unido en la creación (Col 1:17) y lo que se reúne al final (Ef 1:10). ``Sincronía'' (\grc{σύν} + \grc{χρόνος}) es tiempo compartido: eso hace la sabiduría.
\end{column}
\begin{column}{0.43\textwidth}
\precision{Col 1:15--17 usa lenguaje de la Sabiduría (Pr 8:22; Sab 7:26): Pablo atribuye a Cristo lo que Israel decía de la Sabiduría (Dunn, 1998; Wright, 2013). En la LXX de Job 28:23, Dios ``bien \emph{compuso}'' (\grc{συνέστησεν}) el camino de la sabiduría.}
\end{column}
\end{columns}
\end{frame}

% ------------------------------------------------------------------
\section{Vivirlo}

\begin{frame}{Cómo se adquiere la sabiduría}
\scriptsize\renewcommand{\arraystretch}{1.18}
\begin{tabular}{@{}>{\RaggedRight\arraybackslash}p{2.4cm}>{\RaggedRight\arraybackslash}p{1.9cm}>{\RaggedRight\arraybackslash}p{7.5cm}@{}}
\toprule
\textbf{Camino} & \textbf{Texto} & \textbf{RVR1960} \\ \midrule
El temor del Señor & Pr 9:10 & ``El temor de \SENOR{} es el principio de la sabiduría, y el conocimiento del Santísimo es la inteligencia.'' \\
Pedirla & Stg 1:5 & ``Y si alguno de vosotros tiene falta de sabiduría, pídala a Dios, el cual da a todos abundantemente y sin reproche, y le será dada.'' \\
La Palabra & Sal 19:7 [19:8] & ``El testimonio de \SENOR{} es fiel, que hace sabio al sencillo.'' \\
 & Sal 119:98 & ``Me has hecho más sabio que mis enemigos con tus mandamientos, porque siempre están conmigo.'' \\
 & Sal 119:130 & ``La exposición de tus palabras alumbra; hace entender a los simples.'' \\
 & 2~Ti 3:15 & ``las Sagradas Escrituras, las cuales te pueden hacer sabio para la salvación por la fe que es en Cristo Jesús.'' \\
Contar los días & Sal 90:12 & ``Enséñanos de tal modo a contar nuestros días, que traigamos al corazón sabiduría.'' \\
\bottomrule
\end{tabular}\par\smallskip
\precision{\heb{מַחְכִּימַת פֶּתִי} (Sal 19:8): ``que hace sabio al ingenuo''. El \heb{פֶּתִי} (\emph{petí}) es el que se deja llevar por cualquiera.}
\end{frame}

\begin{frame}{Dos sabidurías}
\begin{columns}[t,onlytextwidth]
\begin{column}{0.49\textwidth}
\begin{cita}porque esta sabiduría no es la que desciende de lo alto, sino terrenal, animal, diabólica.\fuente{(Stg 3:15, RVR1960)}\end{cita}
\scriptsize\grc{ἐπίγειος, ψυχική, δαιμονιώδης}: de la tierra, de lo meramente natural, de los demonios. Produce ``celos y contención'' (v. 16).\par\medskip
\begin{cita}Profesando ser sabios, se hicieron necios.\fuente{(Ro 1:22, RVR1960)}\end{cita}
\end{column}
\begin{column}{0.49\textwidth}
\begin{cita}Pero la sabiduría que es de lo alto es primeramente pura, después pacífica, amable, benigna, llena de misericordia y de buenos frutos, sin incertidumbre ni hipocresía.\fuente{(Stg 3:17, RVR1960)}\end{cita}
\scriptsize Siete rasgos: la sabiduría de lo alto se reconoce por su carácter, no por su brillo. Es la sabiduría \destacar{santa} de Charnock: pura antes que eficaz.
\end{column}
\end{columns}
\end{frame}

\begin{frame}{¿Quién la da a conocer?}
\begin{columns}[t,onlytextwidth]
\begin{column}{0.50\textwidth}\small
\begin{itemize}
\item La \destacar{iglesia}: ``para que la multiforme sabiduría de Dios sea ahora dada a conocer por medio de la iglesia'' (Ef 3:10).
\item La \destacar{enseñanza mutua}: ``La palabra de Cristo more en abundancia en vosotros, enseñándoos y exhortándoos unos a otros en toda sabiduría'' (Col 3:16).
\item La \destacar{predicación}: ``a quien anunciamos\ldots{} enseñando a todo hombre en toda sabiduría'' (Col 1:28).
\item La \destacar{oración}: ``os dé espíritu de sabiduría y de revelación'' (Ef 1:17).
\end{itemize}
\end{column}
\begin{column}{0.48\textwidth}
\begin{cita}El Cordero que fue inmolado es digno de tomar el poder, las riquezas, la sabiduría, la fortaleza, la honra, la gloria y la alabanza.\fuente{(Ap 5:12, RVR1960)}\end{cita}
\small La historia termina con la sabiduría devuelta a quien le pertenece: al Cordero.
\end{column}
\end{columns}
\end{frame}

\begin{frame}{Cuando duele}
\begin{columns}[t,onlytextwidth]
\begin{column}{0.50\textwidth}\small
\begin{itemize}
\item La sabiduría de Dios no se ve siempre \emph{en} el momento; se ve \emph{desde} el final (Gn 50:20).
\item José esperó trece años; Pablo escribió Ro 8:28 a creyentes que sufrían ``tribulación, angustia, persecución'' (8:35).
\item El que no entiende puede confiar, porque conoce al que entiende: ``Dios entiende el camino de ella'' (Job 28:23).
\end{itemize}
\end{column}
\begin{column}{0.48\textwidth}
\begin{cita}Fíate de \SENOR{} de todo tu corazón, y no te apoyes en tu propia prudencia. Reconócelo en todos tus caminos, y él enderezará tus veredas.\fuente{(Pr 3:5--6, RVR1960)}\end{cita}
\scriptsize \heb{וְאֶל־בִּינָתְךָ אַל־תִּשָּׁעֵן}: ``y en tu \emph{biná} no te apoyes''. La sabiduría humana es buena; no es un fundamento.
\end{column}
\end{columns}
\end{frame}

% ------------------------------------------------------------------
\section{Sincronía}

\begin{frame}{Reflexión final: la sincronía perfecta}
\begin{columns}[c,onlytextwidth]
\begin{column}{0.56\textwidth}
\includegraphics[width=\linewidth]{bandada_2.jpg}\par\vspace{1pt}
{\fontsize{5}{6}\selectfont\color{gris}\textit{Nota.} Elaboración propia: simulación de una bandada de 5000 aves con el modelo de Reynolds (1987); cada ave atiende solo a sus 7 vecinos más cercanos, como observaron Ballerini et al. (2008) en estorninos.\par}
\end{column}
\begin{column}{0.42\textwidth}\small
Miles de estorninos giran al atardecer como un solo cuerpo. Ninguno dirige; cada uno atiende a seis o siete vecinos; nadie choca.\par\medskip
Las abejas van y vuelven, los salmones remontan el río en que nacieron, el río baja al mar y el mar no se llena (Ec 1:7). Cada movimiento ocurre a su tiempo, uno al lado del otro, sin hacerse daño.\par\medskip
\destacar{Es una sincronía perfecta.}
\end{column}
\end{columns}
\end{frame}

\begin{frame}{La naturaleza que se orquesta}
\fontsize{6.6}{7.8}\selectfont\renewcommand{\arraystretch}{1.12}
\begin{tabular}{@{}>{\RaggedRight\arraybackslash}p{2.0cm}>{\RaggedRight\arraybackslash}p{5.2cm}>{\RaggedRight\arraybackslash}p{4.5cm}@{}}
\toprule
\textbf{Fenómeno} & \textbf{Lo que se observa} & \textbf{Eco bíblico} \\ \midrule
Estorninos & bandadas de miles giran juntas; cada ave sigue a 6--7 vecinos (Ballerini et al., 2008) & las langostas ``no tienen rey, y salen todas por cuadrillas'' (Pr 30:27) \\
Luciérnagas & en algunos bosques miles encienden su luz al mismo tiempo, sin un director (Strogatz, 2003) & ``Todo lo hizo hermoso en su tiempo'' (Ec 3:11) \\
Abejas & la danza de la obrera indica a la colmena la dirección y la distancia del alimento (von Frisch, 1967) & ``la hormiga\ldots{} en el verano preparan su comida'' (Pr 30:25) \\
Salmones & vuelven a desovar al río donde nacieron, guiados por el olor de sus aguas (Hasler \& Scholz, 1983) & ``la cigüeña en el cielo conoce su tiempo'' (Jer 8:7) \\
El agua & los ríos van al mar, el mar no se llena, y el agua vuelve a correr & Ec 1:7; ``destilan rocío los cielos'' (Pr 3:20) \\
\bottomrule
\end{tabular}\par\smallskip
\begin{cita}[fontupper=\scriptsize\itshape,before skip=2pt,after skip=2pt]Cuatro cosas son de las más pequeñas de la tierra, y las mismas son más sabias que los sabios.\fuente{(Pr 30:24, RVR1960)}\end{cita}
\scriptsize\heb{חֲכָמִים מְחֻכָּמִים} (\emph{jajamim mejukkamim}): ``sabias, hechas sabias''. Su sabiduría es recibida. La LXX dice de las langostas que marchan ``a una sola orden, en buen orden'' (\grc{ἀφ᾽ ἑνὸς κελεύσματος εὐτάκτως}).
\end{frame}

\begin{frame}{Para ver la sincronía}
\begin{columns}[c,onlytextwidth]
\begin{column}{0.56\textwidth}
\includegraphics[width=0.49\linewidth]{bandada_1.jpg}\hfill\includegraphics[width=0.49\linewidth]{bandada_3.jpg}\par\vspace{1pt}
{\fontsize{5}{6}\selectfont\color{gris}\textit{Nota.} Elaboración propia: dos momentos de la misma simulación.\par}\medskip
\small Video sugerido: \emph{Murmuration}, de Sophie Windsor Clive y Liberty Smith (2011): una bandada de estorninos sobre el río Shannon, Irlanda, filmada desde una canoa.
\end{column}
\begin{column}{0.42\textwidth}
\centering
\qrcode[height=2.6cm]{https://vimeo.com/31158841}\par\smallskip
\footnotesize\href{https://vimeo.com/31158841}{vimeo.com/31158841}\par\medskip
\begin{cita}Aun la cigüeña en el cielo conoce su tiempo, y la tórtola y la grulla y la golondrina guardan el tiempo de su venida; pero mi pueblo no conoce el juicio de \SENOR.\fuente{(Jer 8:7, RVR1960)}\end{cita}
\end{column}
\end{columns}
\precision{\heb{חֲסִידָה} (\emph{jasidá}), ``cigüeña'', tiene la raíz de \emph{jésed}. Las aves guardan sus tiempos (\heb{מוֹעֲדֶיהָ}); el pueblo olvida el de Dios.}
\end{frame}

\begin{frame}{Pero la sabiduría es más que la naturaleza}
\begin{columns}[t,onlytextwidth]
\begin{column}{0.48\textwidth}
\begin{cita}[fontupper=\scriptsize\itshape]Los cielos cuentan la gloria de Dios, y el firmamento anuncia la obra de sus manos. Un día emite palabra a otro día, y una noche a otra noche declara sabiduría.\fuente{(Sal 19:1--2 [19:2--3], RVR1960)}\end{cita}
\small La creación \emph{cuenta} la gloria. Pero la sabiduría de Dios se ve todavía más en:
\end{column}
\begin{column}{0.50\textwidth}
\centering
\begin{tikzpicture}[font=\scriptsize,baseline=(current bounding box.north)]
 \node[caja,text width=4.4cm,align=left] (a) at (0,1.95) {\textbf{la creación del ser humano}: formado, tejido, bordado (Sal 139)};
 \node[caja,text width=4.4cm,align=left] (b) at (0,0.95) {\textbf{la soberanía}: quita reyes y pone reyes (Dn 2:21)};
 \node[caja,text width=4.4cm,align=left] (c) at (0,-0.2) {\textbf{el mal convertido en bien}: José (Gn 50:20)};
 \node[draw=dorado,line width=1.4pt,fill=azul,text=papiro,rounded corners=3pt,text width=4.4cm,align=left,inner sep=5pt] (d) at (0,-1.45) {\textbf{el plan de redención}: Cristo crucificado, sabiduría de Dios (1~Co 1:24); la iglesia de judíos y gentiles (Ef 3:10)};
\end{tikzpicture}
\end{column}
\end{columns}
\end{frame}

\begin{frame}{El sincronizador de los atributos}
\begin{columns}[c,onlytextwidth]
\begin{column}{0.56\textwidth}
\centering
\begin{tikzpicture}[font=\scriptsize]
 \gear{(0,0)}{1.05}{18}{azul}
 \fill[papiro] (0,0) circle (0.78);
 \node[align=center,text=azul,font=\scriptsize\bfseries] at (0,0) {Sabiduría\\ \heb{חָכְמָה}};
 \foreach \a/\t in {90/Santidad,30/Justicia,-30/Misericordia,-90/Fidelidad,-150/Amor,150/Poder}{
   \gear{(\a:2.35)}{0.62}{12}{dorado}
   \fill[papiro] (\a:2.35) circle (0.47);
   \node[font=\tiny\bfseries,text=azuloscuro] at (\a:2.35) {\t};}
\end{tikzpicture}\par\vspace{1pt}
{\fontsize{5}{6}\selectfont\color{gris}\textit{Nota.} Elaboración propia a partir de la imagen ``El mecanismo divino'' del material del autor.\par}
\end{column}
\begin{column}{0.42\textwidth}\small
La sabiduría es \destacar{superlativa}: el sincronizador perfecto de todos los atributos. Por ella el poder nunca actúa sin amor, la justicia nunca sin misericordia, la misericordia nunca contra la santidad.\par\medskip
Junto con la \destacar{santidad}, es un atributo que se ve en todos los demás: la santidad es su \emph{brillo}; la sabiduría, su \emph{concierto}.\par\medskip
\precision{lectura propia de la serie. La imagen del engranaje es una analogía: en Dios no hay piezas (Charnock, 1682/1996).}
\end{column}
\end{columns}
\end{frame}

\begin{frame}{Definición: la sabiduría como sincronía}
\begin{tcolorbox}[colback=papiro,colframe=azul,boxrule=1pt,arc=3pt]
\small La sabiduría de Dios es la perfección eterna por la cual Él \destacar{concibe}, \destacar{ordena} y \destacar{conduce} todas las cosas ---desde el átomo hasta la historia de la redención--- hacia su fin, de modo que cada parte se mueve en su tiempo y en su lugar, en concierto con todas las demás, sin choque ni desperdicio; y todos sus atributos obran en Él como una sola voluntad: reunir todas las cosas en Cristo para su gloria (Pr 3:19; Ro 11:36; Ef 1:10; Col 1:17).
\end{tcolorbox}
\centering
\begin{tikzpicture}
 \foreach \x/\n/\t/\d/\r in {0/1/{La partitura}/{el diseño eterno, escrito antes de que existieran tiempo, espacio y materia}/{Pr 8:22--23; 1~Co 2:7}, 4.2/2/{La batuta}/{la providencia que marca a cada criatura su entrada y su tiempo}/{Dn 2:21; Ro 8:28}, 8.4/3/{El acorde final}/{la redención: todo reunido bajo una cabeza, Cristo}/{Ef 1:10; Ap 5:12}}{
  \node[caja,anchor=north] (b\n) at (\x,0) {\parbox[t][1.55cm][t]{3.55cm}{\RaggedRight
    {\small\textbf{\textcolor{azul}{\t}}}\\[2pt]\scriptsize \d\par\vfill{\tiny\color{gris}(\r)}}};
  \node[circle,fill=azul,text=papiro,font=\scriptsize\bfseries,inner sep=1.6pt] at (b\n.north west) {\n};}
\end{tikzpicture}\par\smallskip
\small Si el universo fuera una orquesta, la sabiduría de Dios no sería un instrumento más: sería la partitura y la batuta a la vez.
\end{frame}

% ------------------------------------------------------------------
\section{Síntesis}

\begin{frame}{Síntesis}
\begin{tcolorbox}[colback=papiro,colframe=azul,boxrule=1pt,arc=3pt]
\small La \emph{jojmá} de Dios es destreza perfecta: con ella fundó la tierra (Pr 3:19), estuvo ``ordenándolo todo'' (Pr 8:30), formó al ser humano (Sal 139), gobierna los tiempos (Dn 2:21) y convirtió en bien el mal hecho a José (Gn 50:20). Israel la bendice cada día; Pablo la ve en Cristo crucificado, ``poder de Dios y sabiduría de Dios'' (1~Co 1:24), en quien todas las cosas subsisten (Col 1:17) y serán reunidas (Ef 1:10). Quien la recibe deja de gloriarse en su propia sabiduría y se gloría en el Señor (Jer 9:23--24; 1~Co 1:31).
\end{tcolorbox}
\scriptsize\renewcommand{\arraystretch}{1.15}
\begin{tabular}{@{}>{\RaggedRight\arraybackslash}p{2.4cm}>{\RaggedRight\arraybackslash}p{9.4cm}@{}}
\toprule
\textbf{Mirada} & \textbf{Aporte} \\ \midrule
Hebrea & \emph{jojmá} como destreza; Pr 8 (\emph{amón}); Pr 3:19--20; Bezalel, Salomón, Jetro; José (\emph{jashav}). \\
Judía & la Torá como plano (\emph{Bereshit Rabá} 1:1); Rashi; ``el que ve lo que va a nacer''; \emph{Asher yatsar}; Sabiduría y Ben Sirá; 1QS y 11QSal\textsuperscript{a}. \\
Pink & los decretos son sabios: mejores fines, medios más adecuados; la máquina; sabiduría santa. \\
Paulina & Cristo, sabiduría de Dios; Jer 9:24 en 1~Co 1:31; Ro 11:33--36; Ef 3:10; Col 1:17; Ef 1:10. \\
\bottomrule
\end{tabular}
\end{frame}

\begin{frame}{Para la reflexión}
\begin{cita}Mas alábese en esto el que se hubiere de alabar: en entenderme y conocerme, que yo soy \SENOR, que hago misericordia, juicio y justicia en la tierra; porque estas cosas quiero, dice \SENOR.\fuente{(Jer 9:24 [9:23], RVR1960)}\end{cita}
\begin{center}\heb{הַשְׂכֵּל וְיָדֹעַ אוֹתִי}\end{center}
\small
\begin{itemize}
\item ¿En qué sabiduría propia me estoy gloriando?
\item ¿Qué parte de mi historia todavía no entiendo, y puedo confiarla al que ``entiende el camino''?
\item ¿Dónde veo hoy, en lo pequeño, la sincronía de Dios?
\item ¿Le he pedido sabiduría esta semana (Stg 1:5)?
\end{itemize}
\end{frame}

% ------------------------------------------------------------------
\section{Glosario}

\begin{frame}{Glosario}
\scriptsize\RaggedRight
\begin{columns}[t,onlytextwidth]
\begin{column}{0.49\textwidth}

\par\noindent\hangindent=1em{\color{azul}\textbf{\emph{Amón}}} (\heb{אָמוֹן}) artesano, maestro de obras; o niño criado (Pr 8:30).\par\vspace{2pt}
\par\noindent\hangindent=1em{\color{azul}\textbf{\emph{Anakephalaiōsis}}} (\grc{ἀνακεφαλαίωσις}) recapitulación: reunir todo bajo una cabeza (Ef 1:10).\par\vspace{2pt}
\par\noindent\hangindent=1em{\color{azul}\textbf{\emph{Biná}}} (\heb{בִּינָה}) discernimiento, entendimiento.\par\vspace{2pt}
\par\noindent\hangindent=1em{\color{azul}\textbf{\emph{Da'at}}} (\heb{דַּעַת}) conocimiento por experiencia.\par\vspace{2pt}
\par\noindent\hangindent=1em{\color{azul}\textbf{\emph{Etsá}}} (\heb{עֵצָה}) consejo, designio.\par\vspace{2pt}
\par\noindent\hangindent=1em{\color{azul}\textbf{\emph{Jajam}}} (\heb{חָכָם}) sabio, hábil.\par\vspace{2pt}
\par\noindent\hangindent=1em{\color{azul}\textbf{\emph{Jashav}}} (\heb{חָשַׁב}) pensar, planear, diseñar (Gn 50:20; Éx 31:4).\par\vspace{2pt}
\par\noindent\hangindent=1em{\color{azul}\textbf{\emph{Jojmá}}} (\heb{חָכְמָה}) sabiduría, destreza.\par\vspace{2pt}

\end{column}
\begin{column}{0.49\textwidth}

\par\noindent\hangindent=1em{\color{azul}\textbf{\emph{Majashavá}}} (\heb{מַחֲשָׁבָה}) diseño, plan, pensamiento.\par\vspace{2pt}
\par\noindent\hangindent=1em{\color{azul}\textbf{\emph{Mystērion}}} (\grc{μυστήριον}) plan oculto y revelado (1~Co 2:7).\par\vspace{2pt}
\par\noindent\hangindent=1em{\color{azul}\textbf{\emph{Polypoíkilos}}} (\grc{πολυποίκιλος}) multiforme, de muchos colores (Ef 3:10).\par\vspace{2pt}
\par\noindent\hangindent=1em{\color{azul}\textbf{\emph{Sophía}}} (\grc{σοφία}) sabiduría.\par\vspace{2pt}
\par\noindent\hangindent=1em{\color{azul}\textbf{\emph{Synergéō}}} (\grc{συνεργέω}) cooperar (Ro 8:28).\par\vspace{2pt}
\par\noindent\hangindent=1em{\color{azul}\textbf{\emph{Synístēmi}}} (\grc{συνίστημι}) mantener unido, subsistir (Col 1:17).\par\vspace{2pt}
\par\noindent\hangindent=1em{\color{azul}\textbf{\emph{Tevuná}}} (\heb{תְּבוּנָה}) inteligencia, entendimiento.\par\vspace{2pt}
\par\noindent\hangindent=1em{\color{azul}\textbf{\emph{Tushiyá}}} (\heb{תּוּשִׁיָּה}) sabiduría eficaz (Is 28:29).\par\vspace{2pt}

\end{column}
\end{columns}
\end{frame}

% ------------------------------------------------------------------
\section{Referencias}

\begin{frame}{Referencias}
\fontsize{5.4}{6.4}\selectfont
\begin{columns}[t,onlytextwidth]
\begin{column}{0.49\textwidth}
\refapa{Ballerini, M., Cabibbo, N., Candelier, R., Cavagna, A., Cisbani, E., Giardina, I., Lecomte, V., Orlandi, A., Parisi, G., Procaccini, A., Viale, M., \& Zdravkovic, V. (2008). Interaction ruling animal collective behavior depends on topological rather than metric distance: Evidence from a field study. \emph{Proceedings of the National Academy of Sciences, 105}(4), 1232--1237. https://doi.org/10.1073/pnas.0711437105}
\refapa{Bauer, W., Danker, F.~W., Arndt, W.~F., \& Gingrich, F.~W. (2000). \emph{A Greek-English lexicon of the New Testament and other early Christian literature} (3.ª ed.). University of Chicago Press.}
\refapa{Berkhof, L. (1938). \emph{Systematic theology}. Eerdmans.}
\refapa{Charnock, S. (1996). \emph{The existence and attributes of God} (Vols. 1--2). Baker Books. (Obra original publicada en 1682)}
\refapa{Dunn, J.~D.~G. (1998). \emph{The theology of Paul the apostle}. Eerdmans.}
\refapa{Elbogen, I. (1993). \emph{Jewish liturgy: A comprehensive history} (R.~P. Scheindlin, Trad.). Jewish Publication Society.}
\refapa{Elliger, K., \& Rudolph, W. (Eds.). (1997). \emph{Biblia Hebraica Stuttgartensia} (5.ª ed.). Deutsche Bibelgesellschaft.}
\refapa{Hasler, A.~D., \& Scholz, A.~T. (1983). \emph{Olfactory imprinting and homing in salmon}. Springer.}
\refapa{Koehler, L., \& Baumgartner, W. (2001). \emph{The Hebrew and Aramaic lexicon of the Old Testament} (M.~E.~J. Richardson, Trad.; Ed. de estudio). Brill.}
\refapa{McNamara, M. (1992). \emph{Targum Neofiti 1: Genesis} (The Aramaic Bible, Vol. 1A). Liturgical Press.}
\refapa{Nestle, E., \& Aland, K. (Eds.). (2012). \emph{Novum Testamentum Graece} (28.ª ed.). Deutsche Bibelgesellschaft.}
\end{column}
\begin{column}{0.49\textwidth}
\refapa{Pink, A.~W. (2020). \emph{Los atributos de Dios} (E. Rodríguez Cruz, Trad.). Chapel Library. (Obra original publicada en 1930)}
\refapa{Qumran-Digital. (2022). \emph{1QS (transcripción)}. Akademie der Wissenschaften zu Göttingen. https://lexicon.qumran-digital.org}
\refapa{Rahlfs, A., \& Hanhart, R. (Eds.). (2006). \emph{Septuaginta} (Ed. rev.). Deutsche Bibelgesellschaft.}
\refapa{Rashi. (s.f.). \emph{Comentario a Éxodo y Proverbios}. Sefaria. https://www.sefaria.org}
\refapa{Reynolds, C.~W. (1987). Flocks, herds and schools: A distributed behavioral model. \emph{Computer Graphics, 21}(4), 25--34.}
\refapa{Sanders, J.~A. (1965). \emph{The Psalms scroll of Qumrân Cave 11 (11QPs\textsuperscript{a})} (Discoveries in the Judaean Desert IV). Clarendon Press.}
\refapa{\emph{Santa Biblia: Reina-Valera 1960}. (1960). Sociedades Bíblicas en América Latina.}
\refapa{Schreier, J. (Trad.). (2022). \emph{The Sefaria Midrash Rabbah} [Bereshit Rabá 1:1]. Sefaria. https://www.sefaria.org}
\refapa{Steinsaltz, A. (Ed.). (2017). \emph{Koren Talmud Bavli: Edición William Davidson} [Berajot 58a, 60b; Tamid 32a]. Sefaria. https://www.sefaria.org}
\refapa{Strogatz, S. (2003). \emph{Sync: The emerging science of spontaneous order}. Hyperion.}
\refapa{von Frisch, K. (1967). \emph{The dance language and orientation of bees}. Harvard University Press.}
\refapa{Windsor Clive, S., \& Smith, L. (2011). \emph{Murmuration} [Video]. Vimeo. https://vimeo.com/31158841}
\refapa{Wright, N.~T. (2013). \emph{Paul and the faithfulness of God}. Fortress Press.}
\end{column}
\end{columns}
\par\smallskip{\color{gris}Las citas bíblicas son de la Reina-Valera 1960; ``Jehová'' se reemplaza por \SENOR. Hebreo: BHS; griego: NA28 y LXX de Rahlfs-Hanhart. Fuentes judías: Sefaria; lo rabínico se marca como \emph{precisión}. \emph{Pirké Avot} 4:1 se cita de la Misná (Sefaria).}
\end{frame}

\end{document}
